\addcontentsline{toc}{section}{ВВЕДЕНИЕ}
\section*{ВВЕДЕНИЕ}
Случайные графы находят широкое приложение в различных областях,
таких как физика, биология, социальные и экономические науки. 
Наиболее теоретически исследованной моделью случайного графа является 
модель Эрдёша--Реньи $G(n,p)$, в которой строится 
неориентированный граф на $n$ вершинах, каждое ребро добавляется с
вероятностью  $p$ независимо от остальных. Данная модель имеет 
важное свойство -- фазовый переход, который заключается в резком изменении структуры графа, после достижения критического значения параметра $p$. 

Несмотря на то, что фазовый переход в модели $G(n,p)$ 
теоретически хорошо изучен, её поведение  
при конечных $n$ существенно отличается от асимптотического. Это делает численное исследование 
модели актуальным для проверки скорости сходимости теоретических
результатов.

Целью работы является численное исследование модели Эрдёша--Реньи. 
Были реализованы графы, алгоритмы на графах и вычислительный эксперимент. Также были визуализированы результаты экспериментов.
Были изучены характеристики полученных графов -- количество 
компонент связности, размер наибольшей компоненты связности, структура подграфа, построенного на вершинах из компоненты связности, размер наибольшего найденного пути, асимметричность распределения размеров наибольших компонент связности.
Для достижения цели решались следующие задачи:
\begin{enumerate}
    \item реализация генератора случайных графов 
        согласно модели $G(n,p)$ на языке программирования C++;
    \item реализация алгоритмов для вычисления характеристик графов;
    \item проведение вычислительных экспериментов;
    \item анализ и визуализация полученных данных;
\end{enumerate}

Работа состоит из введения, трёх глав, заключения и списка 
использованных источников. В первой главе рассматриваются 
основные понятия теории графов и свойства модели $G(n,p)$.
Во второй главе представлена программная реализация графов, 
вычислительного эксперимента и анализа полученных результатов.
В третьей главе приводятся результаты экспериментов.
